\documentclass[a4paper,12pt]{article}
\usepackage[T2A]{fontenc}
\usepackage[utf8]{inputenc}
\usepackage[russian]{babel}
\usepackage{amsmath,amssymb,mathtools}
\usepackage{geometry}
\geometry{left=24mm,right=22mm,top=24mm,bottom=23mm}
\usepackage{helvet}
\renewcommand{\familydefault}{\sfdefault}
\usepackage{microtype,booktabs,tabularx,enumitem,xcolor,tikz,fancyhdr,hyperref}
\definecolor{accent}{HTML}{008080}
\definecolor{darkaccent}{HTML}{004D4D}
\definecolor{graytext}{HTML}{42484A}
\definecolor{error}{HTML}{A33333}
\hypersetup{colorlinks=true,linkcolor=darkaccent,urlcolor=darkaccent}
\setlength{\parindent}{0pt}
\setlength{\parskip}{6pt}
\linespread{1.08}
\setlist[itemize]{leftmargin=*,itemsep=2pt,topsep=3pt}
\setlength{\headheight}{15pt}
\pagestyle{fancy}
\fancyhf{}
\fancyhead[L]{\footnotesize\color{graytext}Физика · Ревью домашней работы}
\fancyhead[R]{\footnotesize\color{graytext}София Лузянина}
\fancyfoot[C]{\footnotesize\color{graytext}\thepage}
\fancyfoot[L]{\scriptsize Лёвин Артём Александрович · эксклюзивно для Софии Лузяниной}
\renewcommand{\headrulewidth}{0.25pt}
\renewcommand{\footrulewidth}{0pt}
\newcommand{\err}[1]{\textcolor{error}{\textbf{Ошибка:} #1}}
\newcommand{\idea}[1]{\textcolor{darkaccent}{\textbf{Ключевая идея:} #1}}
\newcommand{\result}[1]{\textcolor{darkaccent}{\textbf{Итог:} #1}}
\newcommand{\practice}{\medskip\textbf{Задача на закрепление.}~}
\newcommand{\task}[1]{\clearpage\phantomsection\label{task:#1}\section*{Задание №\,#1}}
\begin{document}
\thispagestyle{empty}
\begin{center}
\vspace*{24mm}
{\large\color{darkaccent}ФИЗИКА}\par
\vspace{22mm}
{\LARGE Ревью домашней работы}\par
\vspace{12mm}
{\Large\color{accent}Проекции движения и отражение\par от наклонной плоскости}\par
\vspace{19mm}
{\normalsize Разбор ошибок · Механика}\par
\vfill
{\large София Лузянина}\par
\vspace{10mm}
{\normalsize 23 ноября 2026 г.}\par
\vspace{25mm}
{\small Лёвин Артём Александрович}\par
{\small эксклюзивно для Софии Лузяниной}\par
\vspace{14mm}
\begin{tikzpicture}[x=1cm,y=1cm]
 \draw[accent,line width=0.7pt] (-4,0) .. controls (-2.5,0.75) and (-0.8,-0.8) .. (0,0) .. controls (1.3,0.75) and (2.9,-0.6) .. (4,0);
\end{tikzpicture}
\end{center}
\clearpage
\section*{Сводка проверки}
В итоговое ревью включены задания, в решении которых обнаружена ошибка или существенный некорректный шаг. Полностью верные работы в таблице не перечисляются.
\begin{tabularx}{\linewidth}{@{}p{13mm}X p{37mm}@{}}
\toprule
\textbf{№} & \textbf{Содержание замечания} & \textbf{Результат проверки}\\
\midrule
\hyperref[task:1]{1} & Смешаны горизонтальные и наклонные оси; неверны проекции ускорения & Нужна правка\\
\hyperref[task:2]{2} & Не учтены компоненты тяжести, влияющие на время и дальность & Нужна правка\\
\hyperref[task:3]{3} & Неверное ускорение по нормали; неоднозначно определена высота & Нужна правка\\
\hyperref[task:4]{4} & Применена формула горизонтальной дальности к наклонной поверхности & Нужна правка\\
\hyperref[task:5]{5} & Ошибочно перенесено свойство равных дальностей на наклон & Нужна правка\\
\hyperref[task:8]{8} & Верный ответ при неправильном равенстве вертикальной скорости полной скорости & Частичная правка\\
\bottomrule
\end{tabularx}

\medskip\textbf{Граница проверки.} На предоставленных фотографиях решения заданий №\,6 и №\,7 обрезаны. Восстановить их ход достоверно нельзя; эти задания исключены из ревью. Задания №\,9 и №\,10 проверены и не содержат существенных ошибок.

\medskip\textbf{Уточнение физической модели для заданий №\,1--5.} В формулировке не обозначено, повышается или понижается поверхность в направлении $A\to B$. Далее принято, что линия $AB$ направлена \emph{вниз по склону}: ось $Ox$ идёт от $A$ к $B$, ось $Oy$ направлена по внешней нормали к поверхности. Это согласуется с разрешимым вариантом задания №\,4. Если $AB$ направлена вверх, проекция $a_x$ меняет знак и численные ответы для дальности будут другими. В обоих вариантах $a_y=-g\cos\alpha$.

\medskip\idea{Первый шаг в каждой задаче на наклонной плоскости: обозначить оси, показать направление силы тяжести и получить обе её проекции.}
\clearpage
\begin{center}\vspace*{8cm}{\LARGE\color{darkaccent}Разбор ошибок}\par\vspace{5mm}{\small Задания №\,1--5 и №\,8}\end{center}

\task{1}
\textbf{Условие.} Шайба стартует со скоростью $8$ м/с под углом $\beta=30^\circ$ к линии $AB$. Плоскость наклонена на $\alpha=30^\circ$. Найти $v_{0x}$, $v_{0y}$, $a_x$, $a_y$, $x(t)$ и $y(t)$ для осей вдоль плоскости и по нормали к ней.

\textbf{Ваше решение.} Найдены $v_{0x}=4\sqrt3$ м/с и $v_{0y}=4$ м/с, после чего принято $a_x=0$, $a_y=-5$ м/с$^2$ и записаны $x(t)=4\sqrt3\,t$, $y(t)=4t-2{,}5t^2$.

\err{Для наклонных осей компоненты ускорения нельзя брать как для горизонтальной и вертикальной системы.}

\textbf{Последний правильный шаг:} $v_{0x}=8\cos30^\circ=4\sqrt3$ м/с, $v_{0y}=8\sin30^\circ=4$ м/с.

\textbf{Первый неверный шаг:} $a_x=0$ и $a_y=-g\sin\alpha=-5$ м/с$^2$. Так получилось, потому что на рисунке оси изображены горизонтально и вертикально, хотя условие требует ось $Ox$ вдоль наклона, а $Oy$ перпендикулярно ему.

\idea{Сила тяжести направлена строго вертикально вниз. Для осей вдоль наклона и по нормали её компоненты равны $g\sin\alpha$ вдоль склона и $-g\cos\alpha$ по внешней нормали.}

\textbf{Исправление и полное решение.} При направлении $Ox$ вниз по склону:
\[
\begin{aligned}
 v_{0x}&=8\cos30^\circ=4\sqrt3\ \text{м/с},&
 a_x&=10\sin30^\circ=5\ \text{м/с}^2,\\
 v_{0y}&=8\sin30^\circ=4\ \text{м/с},&
 a_y&=-10\cos30^\circ=-5\sqrt3\ \text{м/с}^2.
\end{aligned}
\]
При $x_0=y_0=0$ уравнения равноускоренного движения имеют вид:
\[
\boxed{x(t)=4\sqrt3\,t+\frac52t^2},\qquad
\boxed{y(t)=4t-\frac{5\sqrt3}{2}t^2}.
\]
\textbf{Проверка знаков.} При движении вниз по склону $a_x>0$; сила тяжести прижимает шайбу к плоскости, поэтому $a_y<0$. Если ось $Ox$ направить вверх по склону, получится $a_x=-5$ и знак квадратичного слагаемого в $x(t)$ поменяется.

\result{Проекции начальной скорости найдены верно. Исправить необходимо обе проекции ускорения и, следовательно, оба уравнения движения.}

\practice Шайба движется с начальной скоростью $10$ м/с под углом $60^\circ$ к линии склона вниз при $\alpha=30^\circ$. Выберите те же оси и запишите $v_{0x}$, $v_{0y}$, $a_x$, $a_y$, $x(t)$, $y(t)$ для $g=10$ м/с$^2$.

\task{2}
\textbf{Условие.} Шайба запущена со скоростью $12$ м/с под углом $45^\circ$ к линии склона; $\alpha=30^\circ$. Найти время возвращения на линию $AB$ и расстояние до точки возвращения.

\textbf{Ваше решение.} Получено $t=12\sqrt2/5$ с, $AB=28{,}8$ м. В расчёте использовано $a_y=-5$ м/с$^2$, а дальность определена как $v_{0x}t$.

\err{Пропущен множитель $\cos\alpha$ в нормальной проекции тяжести и ускорение вдоль склона.}

\textbf{Последний правильный шаг:} $v_{0y}=12\sin45^\circ=6\sqrt2$ м/с и $v_{0x}=6\sqrt2$ м/с.

\textbf{Первый неверный шаг:} $a_y=-5$, тогда как $a_y=-10\cos30^\circ=-5\sqrt3$ м/с$^2$. Кроме того, $x(t)$ содержит слагаемое $a_xt^2/2$.

\idea{Возвращение на поверхность определяется условием $y(t)=0$; расстояние по склону вычисляется по $x(t)$ в найденный ненулевой момент времени.}

\textbf{Полное правильное решение.} Составим уравнение движения по внешней нормали:
\[
y(t)=6\sqrt2\,t-\frac{5\sqrt3}{2}t^2=0.
\]
Вынесем $t$:
\[
t\left(6\sqrt2-\frac{5\sqrt3}{2}t\right)=0.
\]
Корень $t=0$ соответствует моменту старта, поэтому для возвращения берём
\[
\boxed{t=\frac{12\sqrt2}{5\sqrt3}=\frac{4\sqrt6}{5}\ \text{с}}
\approx1{,}96\ \text{с}.
\]
При направлении $Ox$ вниз по склону $a_x=5$ м/с$^2$. Тогда
\[
\begin{aligned}
AB=x(t)&=6\sqrt2\cdot\frac{4\sqrt6}{5}
 +\frac52\left(\frac{4\sqrt6}{5}\right)^2\\
 &=\frac{48\sqrt3}{5}+\frac{48}{5}
 =\boxed{\frac{48(\sqrt3+1)}5\ \text{м}}
 \approx26{,}23\ \text{м}.
\end{aligned}
\]
\textbf{Проверка.} $y(4\sqrt6/5)=0$; время положительное, $x(t)>0$. Для склона, поднимающегося к $B$, получилось бы $48(\sqrt3-1)/5\approx7{,}03$ м.

\result{Время и дальность необходимо пересчитать с учётом обеих проекций ускорения.}

\practice Со скоростью $9$ м/с шайбу запускают под углом $30^\circ$ к склону, спускающемуся к $B$ под углом $30^\circ$. Найдите положительное время возвращения на $AB$ и пройденное вдоль склона расстояние ($g=10$ м/с$^2$).

\task{3}
\textbf{Условие.} $v_0=14$ м/с, $\beta=30^\circ$, $\alpha=30^\circ$. Найти время достижения максимальной координаты по нормали $y$, значение $y_{\max}$ и величину вертикального расстояния относительно линии $AB$.

\textbf{Ваше решение.} Получено $t=1{,}4$ с, $y_{\max}=4{,}9$ м, затем $H=4{,}9\sin30^\circ=2{,}45$ м.

\err{В расчёте максимумa нормальной координаты снова использовано $-g\sin\alpha$ вместо $-g\cos\alpha$; геометрический смысл высоты не разграничен.}

\textbf{Последний правильный шаг:} $v_{0y}=14\sin30^\circ=7$ м/с.

\textbf{Первый неверный шаг:} $a_y=-5$ м/с$^2$. Правильная величина $a_y=-5\sqrt3$ м/с$^2$.

\idea{В вершине траектории относительно нормальной оси скорость $v_y$ равна нулю. Для наклонной поверхности следует различать нормальное расстояние, вертикальную проекцию нормального смещения и вертикальный зазор до наклонной прямой.}

\textbf{Полное правильное решение.}
\[
 v_y(t)=7-5\sqrt3\,t,
 \qquad v_y(t_{\max})=0,
 \qquad\boxed{t_{\max}=\frac7{5\sqrt3}=\frac{7\sqrt3}{15}\ \text{с}}
 \approx0{,}81\ \text{с}.
\]
\[
 y_{\max}=7t_{\max}-\frac{5\sqrt3}{2}t_{\max}^2
 =\frac{7^2}{2\cdot5\sqrt3}
 =\boxed{\frac{49}{10\sqrt3}\ \text{м}}
 \approx2{,}83\ \text{м}.
\]
\textbf{Уточнение к условию.} Наклонная линия $AB$ не имеет единого горизонтального «уровня», поэтому фраза о вертикальной высоте требует определения способа измерения.
\begin{itemize}
\item \textbf{Вертикальная проекция нормального смещения:}
$h_{\mathrm{proj}}=y_{\max}\cos\alpha
=\frac{49}{20}=2{,}45$ м.
\item \textbf{Вертикальный зазор до линии $AB$ при той же горизонтальной координате:}
$h_{\mathrm{gap}}=y_{\max}/\cos\alpha
=\frac{49}{15}\approx3{,}27$ м.
\end{itemize}
Эти расстояния различны, поскольку измерены вдоль разных направлений. В Вашей записи число $2{,}45$ м получилось при неверном $y_{\max}$ и неверном множителе $\sin\alpha$; совпадение с $h_{\mathrm{proj}}$ случайно.

\result{Время максимума и нормальную координату нужно исправить; вертикальную величину обязательно определить геометрически.}

\practice При $v_0=10$ м/с, $\beta=30^\circ$, $\alpha=30^\circ$ вычислите время максимальной координаты по нормали, $y_{\max}$ и вертикальную проекцию этого нормального смещения ($g=10$ м/с$^2$).

\task{4}
\textbf{Условие.} Шайба вылетает со скоростью $10$ м/с над наклонной поверхностью $\alpha=30^\circ$ и возвращается на линию $AB$ на расстоянии $20$ м. Найти угол запуска $0^\circ<\beta<90^\circ$.

\textbf{Ваше решение.} Взято $L=v_0^2\sin(2\beta)/a_y$ и получено $\beta=45^\circ$.

\err{Формула $L=v_0^2\sin(2\beta)/g$ относится к возвращению на горизонтальный уровень при стандартном выборе осей. Для наклонной плоскости она неприменима.}

\textbf{Последний правильный шаг:} Вы верно распознали, что нужно найти угол из условия возвращения в заданную точку.

\textbf{Первый неверный шаг:} подстановка $g\sin\alpha$ в формулу горизонтальной дальности. Такое действие игнорирует ненулевое ускорение вдоль $AB$ и не учитывает $g\cos\alpha$ при нахождении времени.

\idea{Сначала получить время возвращения из нормального движения, затем подставить его в $x(t)$.}

\textbf{Полное правильное решение.} Для склона вниз к $B$:
\[
t=\frac{2v_0\sin\beta}{g\cos\alpha},\qquad
L=v_0\cos\beta\,t+\frac{g\sin\alpha}{2}t^2.
\]
После подстановки и приведения:
\[
L=\frac{2v_0^2\sin\beta\cos(\beta-\alpha)}
{g\cos^2\alpha}.
\]
Применим формулу произведения синусов и косинусов:
\[
L=\frac{v_0^2}{g\cos^2\alpha}
\left[\sin(2\beta-\alpha)+\sin\alpha\right].
\]
При $v_0=10$, $g=10$, $\alpha=30^\circ$, $L=20$:
\[
20=\frac{40}{3}\left[\sin(2\beta-30^\circ)+\frac12\right],
\quad \sin(2\beta-30^\circ)=1.
\]
С учётом $0^\circ<\beta<90^\circ$:
\[
2\beta-30^\circ=90^\circ,
\qquad\boxed{\beta=60^\circ}.
\]
\textbf{Проверка подстановкой.} $t=2$ с, $x=10\cos60^\circ\cdot2+\tfrac52\cdot2^2=10+10=20$ м. При выборе противоположного направления наклона положительное решение для $L=20$ м отсутствует: максимальная дальность тогда составляет $20/3$ м.

\result{Для заданной дальности на склоне вниз правильный угол равен $60^\circ$. Направление наклона должно быть указано в исходной формулировке.}

\practice При той же скорости $10$ м/с и наклоне вниз $30^\circ$ найдите угол $0^\circ<\beta<90^\circ$, обеспечивающий возвращение на расстоянии $40/3$ м по склону.

\task{5}
\textbf{Условие.} Две шайбы стартуют со скоростью $10$ м/с под углами $\beta_1=30^\circ$ и $\beta_2=60^\circ$ относительно склона $\alpha=30^\circ$. Найти времена возвращения и дальности.

\textbf{Ваше решение.} Вы получили $t_1=2$ с и $t_2=2\sqrt3$ с, а также равные дальности $L_1=L_2=10\sqrt3$ м, опираясь на равенство $\sin60^\circ=\sin120^\circ$.

\err{Равенство дальностей для дополнительных углов действует на горизонтальной плоскости, но в общем случае нарушается на наклонной.}

\textbf{Последний правильный шаг:} формулы $v_{0x}=v_0\cos\beta$, $v_{0y}=v_0\sin\beta$ для осей, связанных со склоном.

\textbf{Первый неверный шаг:} использование $a_y=-5$ м/с$^2$ и $x=v_{0x}t$. Для наклонной плоскости $a_y=-5\sqrt3$ м/с$^2$, а $x=v_{0x}t+\tfrac52 t^2$.

\idea{Угол меняет не только нормальную составляющую начальной скорости и время, но и расстояние, пройденное вдоль склона под действием тяжести.}

\textbf{Полное правильное решение.} Время возвращения:
\[
t(\beta)=\frac{2\cdot10\sin\beta}{10\cos30^\circ}
=\frac{4\sin\beta}{\sqrt3}.
\]
Для первого запуска:
\[
 t_1=\frac2{\sqrt3}\ \text{с},\qquad
L_1=10\cos30^\circ\cdot\frac2{\sqrt3}
 +\frac52\cdot\frac43
=10+\frac{10}{3}=\boxed{\frac{40}{3}\ \text{м}}.
\]
Для второго запуска:
\[
 t_2=2\ \text{с},\qquad
L_2=10\cos60^\circ\cdot2+\frac52\cdot4
=10+10=\boxed{20\ \text{м}}.
\]
\textbf{Проверка.} $t_2>t_1$, но $L_2>L_1$. В общем виде
\[
L(\beta)=\frac{v_0^2}{g\cos^2\alpha}
\left[\sin(2\beta-\alpha)+\sin\alpha\right].
\]
В такой формуле сравниваются $\sin(2\beta-\alpha)$, а не $\sin(2\beta)$.

\textbf{Педагогическое замечание к условию.} Фраза «как могут совпадать расстояния» некорректна для указанных данных. На наклонной поверхности они не совпадают. При запуске на горизонтальный уровень два дополнительных угла действительно могут давать равные дальности при разных временах.

\result{При склоне вниз $t_1=2/\sqrt3$ с, $t_2=2$ с; $L_1=40/3$ м, $L_2=20$ м.}

\practice Две шайбы запускают со скоростью $8$ м/с над склоном вниз $30^\circ$ под углами $30^\circ$ и $60^\circ$ к его линии. Найдите оба времени возвращения и обе дальности; сравните результаты ($g=10$ м/с$^2$).

\task{8}
\textbf{Условие.} Шарик падает с высоты $45$ м без начальной скорости, за $2$ с достигает гладкой наклонной плиты под углом $15^\circ$ и отражается упруго. Найти угол вылета и максимальную высоту после удара ($g=10$ м/с$^2$).

\textbf{Ваше решение.} Получены угол вылета $60^\circ$, высота места удара $25$ м и максимальная высота $40$ м. При этом в одном из промежуточных равенств записано, что вертикальная составляющая скорости после отражения равна всей скорости падения $20$ м/с.

\err{После отражения от наклонной плиты сохраняется модуль скорости, однако вертикальная составляющая в общем случае меньше модуля.}

\textbf{Последний правильный шаг:} скорость непосредственно перед ударом $v=gt=20$ м/с; высота точки удара $h_0=45-\tfrac12\cdot10\cdot2^2=25$ м; угол вылета к горизонту $\theta=90^\circ-2\alpha=60^\circ$.

\textbf{Первый неверный шаг:} равенство $v_{y+}=v=20$ м/с. Скорость $20$ м/с является \emph{модулем} после упругого отражения, а вертикальная компонента определяется углом вылета.

\idea{На гладкой неподвижной плите тангенциальная составляющая скорости сохраняется, нормальная меняет знак. Поэтому угол вылета $90^\circ-2\alpha$, а вертикальная компонента $v_{y+}=v\cos(2\alpha)$.}

\textbf{Полное правильное решение.} Перед ударом модуль вертикальной скорости равен
\[
 v=10\cdot2=20\ \text{м/с}.
\]
После отражения при $\alpha=15^\circ$:
\[
 \theta=90^\circ-30^\circ=\boxed{60^\circ},
\]
\[
 v_{x+}=20\sin30^\circ=10\ \text{м/с},\quad
 v_{y+}=20\cos30^\circ=10\sqrt3\ \text{м/с}.
\]
Проверка сохранения модуля скорости:
\[
 \sqrt{v_{x+}^2+v_{y+}^2}=\sqrt{100+300}=20\ \text{м/с}.
\]
Вертикальный подъём после удара:
\[
 \Delta h=\frac{v_{y+}^2}{2g}=
 \frac{(10\sqrt3)^2}{20}=15\ \text{м}.
\]
Максимальная высота над землёй:
\[
 \boxed{H_{\max}=25+15=40\ \text{м}}.
\]
\textbf{Проверка.} В Ваших последующих вычислениях множитель $\sin^2 60^\circ=3/4$ присутствует, поэтому итоговая численная высота верна, несмотря на ошибочную запись вертикальной скорости.

\result{Угол $60^\circ$ и высота $40$ м верны; исправить необходимо запись о вертикальной составляющей после удара.}

\practice Шарик отпускают с высоты $40$ м. Через $2$ с он упруго ударяется о гладкую плиту, наклонённую на $30^\circ$ к горизонту, и отлетает вправо и вверх. Найдите обе составляющие скорости после удара, высоту удара и максимальную высоту шарика над землёй ($g=10$ м/с$^2$).

\clearpage
\section*{Краткая памятка}
\textbf{Движение в системе осей, связанной со склоном.} Если $Ox$ направлена вниз по склону, $Oy$ --- по внешней нормали, то
\[
 a_x=g\sin\alpha,\qquad a_y=-g\cos\alpha.
\]
При возврате на плоскость сначала найдите ненулевой корень $y(t)=0$, затем вычислите $x(t)$.

\textbf{Упругое отражение от гладкой плиты.} При вертикальном падении с модулем скорости $v$ и $0^\circ<\alpha<45^\circ$:
\[
v_{x+}=v\sin(2\alpha),\quad
v_{y+}=v\cos(2\alpha),\quad
\Delta h=\frac{v_{y+}^2}{2g}.
\]
Для $45^\circ<\alpha<90^\circ$ вертикальная составляющая сразу после удара отрицательна, поэтому шарик не совершает дополнительного подъёма над точкой удара.

\vfill
\begin{center}\footnotesize\color{graytext}
Лёвин Артём Александрович · эксклюзивно для Софии Лузяниной
\end{center}
\end{document}
